\cabstract{
    本文围绕基于自然语言描述的 UI 组件代码生成问题，设计并实现了一套面向 Web 场景的代码生成系统。该系统以 Vue 3 单文件组件为目标输出形式，用户输入组件名称、需求描述和约束信息后，系统调用大语言模型生成对应代码，并进一步提供在线预览、历史回看和文件下载能力。本文的研究重点不在于单次代码生成结果本身，而在于构建一条可执行、可追溯、可控制的工程闭环。

在系统方案上，本文采用前后端分离的 B/S 架构。前端负责需求输入、流式结果展示、代码查看和运行预览，后端负责身份认证、项目空间管理、任务创建、模型调用、结构化解析、版本存储和安全控制。针对模型输出不稳定的问题，系统在生成链路中引入格式约束、结构化提取和缺失片段补全机制，以提高结果可用性。针对生成代码可能带来的运行风险，系统设计了后端规则扫描与前端预览拦截相结合的双层安全策略。为了保证系统在实际使用中的稳定性，系统还实现了任务互斥控制、调用限流、异常统一处理和调用日志记录等机制。

在系统实现与测试阶段，本文完成了需求分析、总体设计、数据库设计、核心功能实现以及功能测试、性能测试和安全测试等内容。测试结果表明，系统能够稳定完成任务创建、流式生成、在线预览、版本管理和结果下载等主要业务流程，对越权访问、限流触发和高风险代码输入也能够给出明确处理结果。研究结果说明，自然语言驱动的 UI 组件生成在课程设计和中小规模前端开发场景中具有较好的工程可行性，并可为后续的多轮生成、自动评测和多模型扩展提供基础。
}{自然语言描述，UI组件生成，Vue 3，流式生成，系统实现}
